\chapter{Transjugular Intrahepatic Portosystemic Shunt}

\section*{Overview}
The transjugular intrahepatic portosystemic shunt (TIPS) is a procedure that creates a channel between the portal and hepatic veins to lower portal pressure. The exact approach, setting, and preparation depend on the indication, anatomy, and the patient's health.

\section*{Alternative Names}
TIPS; TIPS procedure; transjugular intrahepatic portosystemic shunt

\section*{Principle}
A stent is placed through the liver to connect a hepatic vein to the portal vein, redirecting portal blood flow and reducing the pressure that drives variceal bleeding and ascites.

\section*{History}
TIPS was developed in the late 1980s and, with the advent of covered stents, became a standard treatment for refractory complications of portal hypertension.

\section*{Why It Is Done}
Performed for variceal bleeding that fails medical or endoscopic therapy, refractory ascites, hepatic hydrothorax, and Budd-Chiari syndrome.

\section*{Is This a Primary or Secondary Test or Procedure}
Secondary (salvage) therapy for complications of portal hypertension.

\section*{How It Is Performed}
Performed under sedation or anesthesia by an interventional radiologist. Through the jugular vein, a needle and catheter cross the liver from a hepatic vein to the portal vein, and a covered stent is deployed to create the shunt.

\section*{Preparation}
Pre-procedure imaging, correction of coagulopathy, and antibiotic prophylaxis.

\section*{Contraindications and Precautions}
Severe heart failure, uncontrolled infection, and advanced liver failure (high MELD) are relative contraindications. Precaution: assess for hepatic encephalopathy risk.

\section*{Risks}
Hepatic encephalopathy, bleeding, infection, shunt stenosis or thrombosis, and heart failure exacerbation.

\section*{Aftercare and Recovery}
Portal pressure is measured; the patient is monitored, and medications target encephalopathy.

\section*{Outcome and Follow-up}
A pressure gradient below 12 mmHg predicts low variceal rebleeding risk.

\section*{Expected Results}
Reduced portosystemic gradient with control of bleeding or ascites.

\section*{Follow-up}
Doppler ultrasound surveillance of shunt patency; interventions for stenosis.

\section*{When to Seek Medical Care}
Follow the procedure team's specific instructions. Seek urgent medical assessment for severe or worsening pain, uncontrolled bleeding, fever or signs of infection, trouble breathing, chest pain, fainting, new weakness or confusion, inability to urinate, or any rapidly worsening symptom. The appropriate warning signs depend on the procedure and the patient's medical condition.

\section*{References}
\begin{enumerate}
  \item MedlinePlus. Transjugular intrahepatic portosystemic shunt (TIPS). U.S. National Library of Medicine; 2025.
  \item Current Medical Diagnosis and Treatment. McGraw-Hill; 2025.
\end{enumerate}

\clearpage
